\section*{Chapter \ref{ch:b3:public-choice-voting} --- Public Choice: Voting}

\begin{solution}{exo:b3:public-choice-voting:1}
From 1870 to 1950 the median rose from 6.5~\% to 17.2~\% of GDP (10.7 points); from 1950 to
2020, to 31.4~\% (14.2 points). The United Kingdom is at the top in 2020, at 48.1~\%, a year
of pandemic spending. % figdata: spending, spending_points
\end{solution}

\begin{solution}{exo:b3:public-choice-voting:2}
25~\%, the median ideal point: against any lower rate, the three voters at 25~\% and above
vote for it; against any higher one, the three at 25~\% and below.
\end{solution}

\begin{solution}{exo:b3:public-choice-voting:3}
$\binom{1000}{500}/2^{1000} = 0.0252$, about $\sqrt{2/(\pi \times 1000)}$: one chance in 40. With millions of
voters and a race that is not exactly even, the probability is far smaller.
\end{solution}

\begin{solution}{exo:b3:public-choice-voting:4}
Platforms that stay apart when the district's voters move; winners whose votes in
office depend on their party rather than on the district's median; or policy that
does not change when the median voter's preferences do.
\end{solution}

\begin{solution}{exo:b3:public-choice-voting:5}
$(1 - 32/40)/0.8 = 25~\%$; with a median of 36{,}000, $(1 - 0.9)/0.8 = 12.5~\%$.
\end{solution}

\begin{solution}{exo:b3:public-choice-voting:6}
With a \omterm{def:b3:public-choice-voting:agenda-setter}{reversion point} of 1 lane, the median household accepts up to 18.5 lanes, and
the setter proposes that. With 6 lanes --- the median's ideal --- every other size is
worse for her, and the setter can win only 6. % figdata: brenmoor_stats
\end{solution}

\begin{solution}{exo:b3:public-choice-voting:7}
$1.02^{20} = 1.486$: the relative price of care has risen by 48.6~\%.
\end{solution}

\begin{solution}{exo:b3:public-choice-voting:8}
The tested solution is \code{ex08_wagner.py}, in the chapter's \code{exercises/}
folder. The slope on log income is 4.91 (s.e.\ 1.04, $R^2$ 0.21, 86 countries) in 2000 and
3.45 (0.88, $R^2$ 0.10, 138 countries) in 2019: weaker, with a larger and different
sample. % output: ex08_wagner
\end{solution}

\begin{solution}{exo:b3:public-choice-voting:9}
No price difference between otherwise identical houses on either side of a
boundary where taxes or services change, or a difference unrelated to the present
value of the tax gap.
\end{solution}

\begin{solution}{exo:b3:public-choice-voting:10}
The proof is in the text: $t = (1 - y_m/\bar y)/a$ falls as $a$ rises. If avoidance raises the
deadweight cost per unit of tax, the median voter faces a higher cost of each point
of tax and chooses less redistribution.
\end{solution}

\begin{solution}{exo:b3:public-choice-voting:11}
Alone, each project gains its district $100 - 40 = 60$ and costs the other two 40 each:
it fails, two to one. If districts 1 and 2 vote for each other's project, both pass:
each of them gains $100 - 80 = 20$, district 3 loses 80. Total benefit is 200 against a
cost of 240: the trade passes an inefficient package.
\end{solution}

\begin{solution}{exo:b3:public-choice-voting:12}
The tested solution is \code{ex12_wagner_within.py}, in the chapter's
\code{exercises/} folder. With country effects (2{,}537 observations, 18 economies),
the slope is 7.55 points of GDP per unit of log income (s.e.\ 1.01); with year effects
added, 1.26 (s.e.\ 4.02), indistinguishable from zero. % output: ex12_wagner_within
Within countries, spending rose as incomes rose, but so did everything else over
time: once the common shocks of each year --- wars, the Depression, the welfare
states after 1945 --- are removed, little of the link remains.
\end{solution}

\begin{solution}{pb:b3:public-choice-voting:pool}
% figdata: brenmoor_stats, brenmoor_u
\textbf{1.} The derivative of $60n\ln q - 20q$ is $60n/q - 20$, zero at $q = 3n$.

\textbf{2.} The second derivative, $-60n/q^2$, is negative: net benefit rises up to $3n$ and
falls after it.

\textbf{3.} 2 members: 550 households have one, 622 two, so households of two or fewer are
58.6~\% and of one 27.5~\%.

\textbf{4.} 6 lanes, the median household's ideal.

\textbf{5.} $\sum_i 60n_i/q = 20N$ gives $q = 3\bar n = 7.35$ lanes: larger households, which value the pool
most, are outvoted, and the median ignores how much they value it.

\textbf{6.} Single people (ideal 3) lose 18.4 dalens a year against their ideal; households
of 3, 4 and 5 lose 13.0, 46.4 and 94.9.

\textbf{7.} $120\ln 2 - 40 = 43.2$.

\textbf{8.} $120\ln q - 20q = 43.2$ above 6 lanes: $q = 13.4$.

\textbf{9.} The median household likes it as much as the old pool, and every household
larger than hers prefers it; with a majority at or above the median, it passes.

\textbf{10.} 18.5 lanes.

\textbf{11.} With no pool, net benefit is zero; the setter proposes the size at which
$120\ln q - 20q = 0$, about 17.0 lanes.

\textbf{12.} Letting voters choose among several sizes, or fixing the \omterm{def:b3:public-choice-voting:agenda-setter}{reversion point}
at a reasonable size rather than at the status quo the setter chooses.

\textbf{13.} Median 23{,}310 dalens, mean 26{,}818.

\textbf{14.} $(1 - 0.869)/0.5 = 26.2~\%$.

\textbf{15.} $0.262 \times 26{,}818 \times (1 - 0.25 \times 0.262) = 6{,}557$ dalens.

\textbf{16.} The decisive voter becomes richer, $y_m/\bar y$ rises, and the chosen rate falls.

\textbf{17.} It cannot settle it. Across countries, more unequal ones spend less, but
their Gini is measured after redistribution, so the arrow can run from spending to
inequality; with the European Gini before transfers there is no relation at all.
The cleaner evidence is within countries: when the franchise widened to poorer
voters, as in question 16 in reverse, welfare spending rose, as the model predicts.

\textbf{18.} $500/0.04 = 12{,}500$ dalens cheaper.

\textbf{19.} The owners at the time the tax gap appears: their house loses 12{,}500 in value;
later buyers pay the tax but a lower price, and are no worse off.

\textbf{20.} With no mobility, nobody needs to move to escape the tax, and prices need
not adjust: the tax falls on residents year by year.

\textbf{21.} A pool that buyers value more than its cost raises Brenmoor's house prices;
one that costs more than it is worth lowers them.

\textbf{22.} Owners pay a tax change once, in the price of their house, the day it is
announced; renters and later buyers do not, so a vote on local taxes is in part a
vote on owners' wealth, and owners have the stronger stake in it.

\textbf{23.} Prices of similar houses that differ across the town boundary by about the
present value of the tax gap, and that move when the gap changes.

\textbf{24.} The referendum chooses a 6-lane pool; the council, against a \omterm{def:b3:public-choice-voting:agenda-setter}{reversion point} of 2
lanes, can win a vote on 13.4 lanes; and a tax 500 dalens a year higher next door is
capitalised as a house-price difference of 12{,}500 dalens at 4~\%.
\end{solution}
